% Part: lambda-calculus
% Chapter: introduction
% Section: basic-pr-lambda

\documentclass[../../../include/open-logic-section]{subfiles}

\begin{document}

\olfileid[pt]{lam}{int}{bas}
\olsection{As funções recursivas primitivas básicas são \usetoken{S}{lambda definable}}

\begin{lem}
As funções $\Zero$, $\Succ$ e $\Proj{n}{i}$ são !!{lambda definable}.
\end{lem}

\begin{proof}
$\Zero$ é simplesmente
$\lambd[x][\lambd[y][y]]$.

A função sucessor~$\Succ$ é definida por
$\fn{Succ}(u) = \lambd[x][\lambd[y][x(uxy)]]$. Você deve pensar por que isso
funciona; para cada numeral $\num{n}$, considerado um iterador, e cada
função~$f$, $\fn{Succ}(\num{n},f)$ é uma função que, com entrada~$y$, aplica
$f$ $n$ vezes a partir de~$y$ e depois a aplica mais uma vez.

Não há nada a dizer sobre projeções: $\fn{Proj}^n_i(x_0, \dots,
x_{n-1}) = x_i$. Em outras palavras, segundo nossas convenções,
$\fn{Proj}^n_i$ é o termo lambda
$\lambd[x_0][\dots \lambd[x_{n-1}][x_i]]$.
\end{proof}

\end{document}

